% Generated draft: requires mathematical and source review.
\documentclass[11pt]{article}
\usepackage[margin=25mm]{geometry}
\usepackage{fontspec}
\setmainfont{Latin Modern Roman}
\usepackage{amsmath,amssymb,mathtools,unicode-math}
\setmathfont{Latin Modern Math}
\usepackage{enumitem,needspace,xurl,hyperref}
\setlength{\parindent}{0pt}
\setlength{\parskip}{5pt}
\setlength{\emergencystretch}{3em}
\newcommand{\N}{\mathbb{N}}
\newcommand{\Z}{\mathbb{Z}}
\newcommand{\Q}{\mathbb{Q}}
\newcommand{\R}{\mathbb{R}}
\newcommand{\C}{\mathbb{C}}
\newcommand{\F}{\mathbb{F}}
\newcommand{\A}{\mathbb{A}}
\newcommand{\PP}{\mathbb{P}}
\newcommand{\E}{\mathbb{E}}
\newcommand{\eps}{\varepsilon}
\newcommand{\dd}{\,\mathrm{d}}
\begin{document}
\section*{UnsolvedMath ID: 30002168}
\textbf{Source problem code:} catalogue entry
\subsection{Definitions and mathematical statement}
The supplied record contains two separate questions, which must not be interpreted as a single conjecture.

First, let
\[
S^2=\{x\in\R^3:\|x\|_2=1\},\qquad
\langle x,y\rangle=\sum_{i=1}^3x_i y_i.
\]
To interpret the numerical bound, provisionally let \(\sigma\) be normalized surface area on \(S^2\), so that \(\sigma(S^2)=1\). This normalization is not explicitly specified in the supplied statement and requires confirmation. Define
\[
\mathcal A=\{A\subseteq S^2:A\text{ is surface measurable and }\langle x,y\rangle\ne0\text{ for every }x,y\in A\},
\qquad
\alpha=\sup_{A\in\mathcal A}\sigma(A).
\]
The orthogonality condition is pointwise for every pair, not merely for almost every pair. The first question asks for the exact value of \(\alpha\) and whether the supplied upper bound \(1/3\) can be improved, meaning whether
\[
\exists\varepsilon>0\quad\forall A\in\mathcal A,
\qquad \sigma(A)\le\frac13-\varepsilon.
\]
The record describes \(1/3\) as a known upper bound; that description is supplied source information, not an independently verified assertion here. The use of a supremum does not assert that an extremizing measurable set exists.

Second, the record asks about the number \(L_n\) of local minima of the Coulomb potential for \(n\) points on \(S^2\). A conventional candidate model, stated explicitly for review, is the collision-free configuration space
\[
X_n=\{(x_1,\ldots,x_n)\in(S^2)^n:x_i\ne x_j\text{ whenever }i\ne j\},\qquad n\ge2,
\]
with energy
\[
E_n(x_1,\ldots,x_n)=\sum_{1\le i<j\le n}\frac1{\|x_i-x_j\|_2}.
\]
Here the distance is the ambient Euclidean distance. Under the non-strict convention, a configuration \(x\in X_n\) is a local minimum if some neighborhood \(U\) of \(x\) in \(X_n\) satisfies \(E_n(x)\le E_n(y)\) for every \(y\in U\). The supplied record does not explicitly fix this energy formula, the configuration-space convention, or the meaning of counting local minima.

In particular, it does not specify whether configurations related by rotations, reflections, or permutations are identified, nor whether strictness is required after any such identification. These choices matter because the conventional energy is invariant under simultaneous rotations. Thus \(L_n\) cannot be assigned an unambiguous mathematical definition from the supplied record alone. Once the intended energy and counting convention are fixed, the requested eventual inequality is
\[
\exists N\in\N\quad\forall n\in\N,\qquad n\ge\max\{N,2\}\Longrightarrow L_n>1.
\]
This second formulation remains conditional on those unresolved conventions.
\subsection{Short English statement}
Two distinct questions are preserved: determine the largest possible normalized area of a measurable subset of \(S^2\) containing no orthogonal pair, including whether the stated bound \(1/3\) can be reduced; and determine whether every sufficiently large number of points admits more than one Coulomb local minimum under the intended counting convention. The area normalization and the definition of distinct local minima require confirmation. Both questions carry the supplied open label, without independent verification.
\subsection{Sources}
\begin{itemize}
\item Local UnsolvedMath registry: \texttt{data/unsolvedmath/problems.json}, with \texttt{catalog/problems.json} as the additional-record registry; supplied permanent ID \(30002168\), exact source code \texttt{catalogue entry}, title ``Orthogonality-Avoiding Sets and Coulomb Local Minima.'' These identify the supplied record; no local file inspection is asserted.
\item Supplied report citation: \emph{Optimal and Near Optimal Configurations on Lattices and Manifolds} (2013), Oberwolfach Reports, pp.\ 2429--2492, DOI \texttt{10.4171/owr/2012/40}. Supplied source URL and literature-source label ``OWR report 40'': \url{https://doi.org/10.4171/owr/2012/40}. The record lists proposed year 2012 and supplies no individual proposer.
\item Supplied literature-source label ``Cohn--Kumar--Schuermann'': \url{https://doi.org/10.1080/10586458.2009.10129052}. No paper title or precise supporting passage is supplied.
\item Literature assessment dated exactly \texttt{2026-08-21}: this is supplied dated curation, not independently verified. It labels the problems open and states that the extracted text conflates the orthogonality-avoiding-set question with the Coulomb local-minima question, with neither the exact extremal value nor the eventual inequality established. The background explicitly requests independent verification of time-sensitive status and cited claims.
\item The supplied statement status is \texttt{corrected\_verified}. Its accompanying verification note says that the official open-problem session places the two questions consecutively and that their separation has been made explicit. This is a supplied verification claim; no source visit or independent verification is asserted here.
\end{itemize}
\end{document}
